\documentclass[12pt]{article}

% Page setup
\usepackage[margin=1in]{geometry}
\usepackage{setspace}
\usepackage{enumitem}
\usepackage{amsmath, amssymb}

% No paragraph indent
\setlength{\parindent}{0pt}

% Line spacing
\renewcommand{\baselinestretch}{1.2}

\begin{document}

\begin{center}
\textbf{\Large MATH 2790 Fall 2025}\\[1em]
\textbf{\Large Homework 6}\\[0.5em]
Due date: Wednesday, Oct 1
\end{center}

\bigskip

\textbf{Instructions:}

\begin{enumerate}[label=(\alph*)]
\item It is expected to show detailed work for every problem. You are allowed to have a
discussion with peers on some challenging questions. But then you should work on
your own to finish the assignment.

\item All questions in this assignment are from sections 3.2--3.3.
\end{enumerate}

\bigskip

\textbf{Problems:}

\begin{enumerate}
\item Question (1) on page 181 of textbook under section 3.2.

\item Question (4) on page 181 of textbook under section 3.2.

\item Question (5) on page 182 of textbook under section 3.2. (Hint: How are PV and the
accumulated value connected? What information can we get from their connection?)

\item Question (1) on page 182 of textbook under section 3.3.

\item Question (2) on page 183 of textbook under section 3.3. (Hint: what is the interest
rate per payment period? Round the interest rate per payment period to 5 decimal
places.)

\item Question (4) on page 183 of textbook under section 3.3. Draw the time diagram first.

\item An annuity has payments of 1000 at the beginning of every quarter for six years.
Another annuity has payments of $X$ at the end of each year for five years. At an
annual effective interest rate of 6\%, the present value of the two annuities are equal.
Find $X$. Round to 4 decimal places.
\end{enumerate}

\end{document}
